% ======================================================================
% Chapter 21: Quantization of Spontaneously Broken Gauge Theories
% (book pages 731-773; PDF pages 754-803)
% ======================================================================
\chapter{Quantization of Spontaneously Broken Gauge Theories}

In Chapter~20 we saw that when a gauge symmetry is spontaneously broken, the
gauge bosons acquire mass.  This phenomenon allowed us to construct a
realistic theory of the weak interactions.  Up to this point, however, we have
discussed spontaneously broken gauge theories only in a simplistic way.  To
isolate the physical degrees of freedom, we have used the device of going to
the \emph{unitarity gauge}.  However, it is not at all clear what the rules of
perturbation theory are in this gauge, or how the unitarity gauge constraint
is maintained when we compute Feynman diagrams.  We have also seen that the
Goldstone bosons that are absorbed into the massive gauge bosons play an
important role in formal arguments about these theories, so we would like to
quantize these theories in a gauge that does not eliminate these particles
from the beginning.

In this chapter we will address these problems, by carrying out the formal
gauge-fixing of theories with spontaneously broken gauge symmetry using the
Faddeev-Popov method.  We will define a class of gauges, called the
\emph{$R_\xi$ gauges}, almost all of which contain the Goldstone bosons of the
original spontaneous symmetry breaking.  These particles cancel the effects of
other unphysical particles in the formalism to maintain the unitarity of the
theory.  These cancellations are a more intricate version of the cancellations
between gauge and ghost degrees of freedom that we saw in Chapter~16.
However, we will see in Section~21.2 that a theory does not forget that it
contains Goldstone bosons and that, under some circumstances, the properties
of the Goldstone bosons in the theory without gauge couplings can carry over
to the theory with massive gauge bosons.

Finally, having defined the perturbation theory and clarified the role of the
Goldstone bosons in spontaneously broken gauge theories, we will carry out
some explicit loop calculations of interest in the theory of weak
interactions.  Here we will see applications of the ideas of Chapter~11, that
a theory with spontaneously broken symmetry can be renormalized with the
counterterms of the symmetric Lagrangian.  In Section~21.3 we will show
through some examples that this result applies with equal force to gauge
theories, and that it endows the weak-interaction gauge theory with
substantial predictive power.

\section{The $R_\xi$ Gauges}

In our discussion of the low-energy effective Lagrangian for weak
interactions, we proposed in Eq.~(20.89) the following expression for
the propagator of a massive gauge boson:
\begin{equation}
  \langle A^\mu(p)A^\nu(-p)\rangle = \frac{i}{p^2 - m^2}\ \Bigl(
    -g^{\mu\nu} + \frac{p^\mu p^\nu}{m^2} \Bigr) .
  \label{eq:21.1}
\end{equation}
This expression is a natural first guess, generalizing the Feynman-'t Hooft
gauge.  However, it is unsatisfactory in a number of ways.

The most important of these defects concerns the treatment of gauge boson
polarization states.  The propagator \eqref{eq:21.1} contains four components,
corresponding to the transverse, longitudinal, and timelike polarizations.
We saw in Chapters~5 and~16 that, for massless gauge bosons, the unphysical
longitudinal and timelike components cancel in computations.  For a massive
gauge boson, however, the longitudinal polarization state corresponds to a
real physical particle; we do not want it to cancel.  Expression \eqref{eq:21.1}
does not take this change into account.

\subsection{An Abelian Example}

To understand this and other formal problems that arise for gauge theories
with spontaneously broken symmetry, we need to carefully redo the
Faddeev-Popov quantization of these theories.  To begin, we will quantize the
spontaneously broken Abelian gauge theory introduced in Eq.~\eqref{eq:20.1}:
\begin{equation}
  \mathcal{L} = -\frac{1}{4}(F_{\mu\nu})^2 + |D_\mu\phi|^2 - V(\phi) ,
  \label{eq:21.2}
\end{equation}
with $D_\mu = \partial_\mu + ieA_\mu$.  Here $\phi(x)$ is a complex scalar
field.  However, it will be most convenient to analyze the model by writing
$\phi$ in terms of its real components,
\begin{equation}
  \phi = \frac{1}{\sqrt{2}}(\phi_1 + i\phi_2) .
  \label{eq:21.3}
\end{equation}
Then the infinitesimal local symmetry transformation is
\begin{equation}
  \delta\phi_1 = -\alpha(x)\phi_2,\qquad \delta\phi_2 = \alpha(x)\phi_1,\qquad
  \delta A_\mu = -\frac{1}{e}\partial_\mu\alpha .
  \label{eq:21.4}
\end{equation}
Let us assume that $V(\phi)$ forces the scalar field to acquire a vacuum
expectation value: $\langle\phi\rangle = v$.  Then we should change variables
by a shift:
\begin{equation}
  \phi_1(x) = v + h(x);\qquad \phi_2 = \varphi .
  \label{eq:21.5}
\end{equation}
The field $\phi_2$ or $\varphi$ is the Goldstone boson.  The Lagrangian
\eqref{eq:21.2} now takes the form
\begin{equation}
  \mathcal{L} = -\frac{1}{4}(F_{\mu\nu})^2
    + \frac{1}{2}\bigl( \partial_\mu h - eA_\mu\varphi \bigr)^2
    + \frac{1}{2}\bigl( \partial_\mu\varphi + eA_\mu(v+h) \bigr)^2
    - V(\phi) .
  \label{eq:21.6}
\end{equation}
This Lagrangian is still invariant under an exact local symmetry,
\begin{equation}
  \delta h = -\alpha\,\varphi,\qquad \delta\varphi = \alpha(v+h),\qquad
  \delta A_\mu = -\frac{1}{e}\partial_\mu\alpha .
  \label{eq:21.7}
\end{equation}
Thus, in order to define the functional integral over the variables
$(\varphi, h, A^\mu)$, we must introduce Faddeev-Popov gauge fixing.

Starting from the functional integral
\begin{equation}
  Z = \int \mathcal{D}A\,\mathcal{D}h\,\mathcal{D}\varphi\ e^{i\int
    d^4x\ \mathcal{L}[A,h,\varphi]} ,
  \label{eq:21.8}
\end{equation}
we can introduce a gauge-fixing constraint as we did in Section~9.4.
Following the steps leading from Eq.~(9.50) to Eq.~\eqref{eq:9.54},
we find
\begin{equation}
  Z = C\cdot \int \mathcal{D}A\,\mathcal{D}h\,\mathcal{D}\varphi\ e^{i\int
    d^4x\ \mathcal{L}[A,h,\varphi]}\ \delta\bigl( G(A,h,\varphi) \bigr)\
    \det\Bigl( \frac{\delta G}{\delta\alpha} \Bigr) ,
  \label{eq:21.9}
\end{equation}
where $C$ is a constant proportional to the volume of the gauge group and
$G(A,h,\varphi)$ is a gauge-fixing condition.  Alternatively, we can introduce
the gauge-fixing constraint as $\delta(G(x) - \omega(x))$ and integrate over
$\omega(x)$ with a Gaussian weight, as in the derivation of Eq.~\eqref{eq:9.56}.
This gives
\begin{equation}
  Z = C'\cdot \int \mathcal{D}A\,\mathcal{D}h\,\mathcal{D}\varphi\ \exp\Bigl[
    i\int d^4x\ \Bigl( \mathcal{L}[A,h,\varphi] - \frac{1}{2\xi}(G)^2
      \Bigr) \Bigr]\ \det\Bigl( \frac{\delta G}{\delta\alpha} \Bigr) .
  \label{eq:21.10}
\end{equation}
The gauge-fixing function $G$ is arbitrary, but we can simplify our formalism
by choosing it appropriately.

An especially convenient choice of the gauge-fixing function is
\begin{equation}
  G = \partial^\mu A_\mu - \xi e v\,\varphi .
  \label{eq:21.11}
\end{equation}
When we form $G^2$, the term quadratic in $A^\mu$ will provide the same
gauge-dependent addition to the gauge field action that we saw in the
derivation of Eqs.~\eqref{eq:9.58} and \eqref{eq:16.29}.  In addition, the
cross term between $A^\mu$ and $\varphi$ is engineered to cancel the quadratic
term of the form $\partial^\mu A_\mu\,\varphi$ coming from the third term of
\eqref{eq:21.6}.  With this choice, the quadratic terms of the gauge-fixed
Lagrangian $(\mathcal{L} - \frac{1}{2\xi}G^2)$ are
\begin{equation}
  \begin{aligned}
    \mathcal{L}_{\mathrm{quad}} &= -\frac{1}{4}(F_{\mu\nu})^2
      - \frac{1}{2\xi}(\partial^\mu A_\mu)^2
      + \frac{1}{2}m_A^2 A_\mu A^\mu + \frac{1}{2}m_A^2\,\varphi^2
      + \frac{1}{2}m_h^2 h^2 .
  \end{aligned}
  \label{eq:21.12}
\end{equation}
The mass term for the $h$ field comes from the expansion of $V(\phi)$, as in
\eqref{eq:20.6}.  The mass term for the gauge field comes from the Higgs
mechanism, that is, from the third term of \eqref{eq:21.6}.  Notice that the
formalism also produces a mass for the Goldstone boson $\varphi$:
\begin{equation}
  m_\varphi^2 = \xi(e v)^2 = \xi m_A^2 .
  \label{eq:21.13}
\end{equation}
The fact that this mass is gauge-dependent is a signal that the Goldstone
boson is a fictitious field, which will not be produced in physical
processes.

To complete the Faddeev-Popov quantization procedure, we must derive the
Lagrangian of the ghosts.  This Lagrangian depends on the gauge variation of
$G$, which can be computed by inserting \eqref{eq:21.7} into \eqref{eq:21.11}.
We find
\begin{equation}
  \delta G = \frac{1}{e}\bigl( -\partial^2 - \xi e^2v^2 \bigr)\alpha
  = \frac{1}{e}\bigl( -\partial^2 - \xi m_A^2 \bigr)\alpha .
  \label{eq:21.14}
\end{equation}
The determinant of this operator can be accounted for by including a set of
Faddeev-Popov ghosts with the Lagrangian,
\begin{equation}
  \mathcal{L}_{\mathrm{ghost}} = \bar{c}\bigl( -\partial^2 - \xi m_A^2 \bigr)c ,
  \label{eq:21.15}
\end{equation}
where $m_A = ev$ as in Eq.~\eqref{eq:21.13}.  Since this is an Abelian gauge
theory, the ghost field does not couple directly to the gauge field.  It
does, however, couple to the physical Higgs field, so it cannot be completely
ignored as in QED.

From the quadratic terms in the Lagrangians for $A^\mu$, $h$, $\varphi$, and
the ghosts, we can readily find the propagators for these fields.  All four
propagators are shown in Fig.~\ref{fig:21.1}.  The only complicated case is
that of the gauge field.  The term in \eqref{eq:21.12} involving $A^\mu$
involves an operator whose Fourier transform is
\begin{equation}
  g^{\mu\nu}k^2 - \Bigl( 1 - \frac{1}{\xi} \Bigr)k^\mu k^\nu
    - m_A^2 g^{\mu\nu} .
  \label{eq:21.16}
\end{equation}
The inverse of this matrix gives the $A^\mu$ field propagator:
\begin{equation}
  \frac{-i}{k^2 - m_A^2}\ \Bigl( g^{\mu\nu} - (1-\xi)\frac{k^\mu k^\nu}
    {k^2 - \xi m_A^2} \Bigr) .
  \label{eq:21.17}
\end{equation}
Notice that the transverse components of the $A$ field and the component $h$
of the Higgs field acquire the masses $m_A$ and $m_h$ that we found in
Section~20.1.  The unphysical components of $A$, the Goldstone bosons, and
the ghosts all acquire the same gauge-dependent mass $\sqrt{\xi}m_A$.

\begin{figure}[t]
  \centering
  \includegraphics[width=0.55\textwidth]{fig21.1.png}
  \caption{Propagators of the gauge field, Higgs fields, and ghosts in the
    Abelian model with spontaneously broken symmetry.}
  \label{fig:21.1}
\end{figure}

\subsection{$\xi$ Dependence in Perturbation Theory}

Because the parameter $\xi$ was introduced only in the gauge fixing, we
expect it to cancel out of all computations of expectation values of
gauge-invariant operators and of $S$-matrix elements.  This cancellation can
be proved to all orders in perturbation theory by using the BRST symmetry of
the gauge-fixed Lagrangian.%
\footnote{See, for example, Taylor (1976).}  Here, however, we will simply
illustrate the cancellation of $\xi$ in a simple example.

\begin{figure}[t]
  \centering
  \includegraphics[width=0.55\textwidth]{fig21.2.png}
  \caption{Diagrams contributing to fermion-fermion scattering at leading
    order in the Abelian model with spontaneous symmetry breaking.}
  \label{fig:21.2}
\end{figure}

Consider coupling a fermion to the spontaneously broken gauge theory through
a chiral interaction:
\begin{equation}
  \mathcal{L}_f = \bar{\psi}_L(i\slashed{D})\psi_L
    + \bar{\psi}_R(i\slashed{\partial})\psi_R
    - \lambda\bigl( \bar{\psi}_R\phi\,\psi_L + \bar{\psi}_L\phi^{*}\psi_R
      \bigr) ,
  \label{eq:21.18}
\end{equation}
with $D_\mu = \partial_\mu + ieA_\mu$ as before.  This is a stripped-down,
Abelian version of the coupling of fermions to the weak interaction gauge
theory.  The fermion $\psi$ receives a mass
\begin{equation}
  m_f = \frac{\lambda v}{\sqrt{2}}
  \label{eq:21.19}
\end{equation}
from the spontaneous symmetry breaking.  (This theory has an axial vector
anomaly that would render loop calculations inconsistent, but we will analyze
it only at the level of tree diagrams.)

In this theory, the leading-order diagrams contributing to fermion-fermion
scattering are those shown in Fig.~\ref{fig:21.2}.  Notice that the
contribution from the exchange of the unphysical particle $\varphi$ must be
included, since this particle appears in the Feynman rules.  The ghosts do
not appear in this process until the one-loop level.  Since the propagator of
the physical Higgs particle $h$ is independent of $\xi$, the cancellation of
the $\xi$ dependence must take place between the transverse and longitudinal
components of $A^\mu$ and the Goldstone boson $\varphi$.

Carrying out this cancellation explicitly, one finds that the contribution
of the longitudinal gauge boson propagator combines with the Goldstone boson
exchange to give a gauge-invariant result.  This verifies that the $R_\xi$
gauges give results independent of $\xi$, as required by gauge invariance.

\subsection{The Feynman-'t Hooft Gauge and the Unitary Gauge}

Among the $R_\xi$ gauges, the choice $\xi = 1$ is particularly convenient.
In this gauge, called the \emph{Feynman-'t Hooft gauge}, the gauge boson and
Goldstone boson propagators take the simple forms
\begin{equation}
  \frac{-ig^{\mu\nu}}{k^2 - m_A^2},\qquad
  \frac{i}{k^2 - m_A^2} .
\end{equation}
This gauge, still called the Feynman-'t Hooft gauge, is the most convenient
one for general higher-order computations.

For any finite value of $\xi$, the gauge boson and Goldstone boson propagators
fall off as $1/k^2$ and thus obey the general power-counting analysis of
Section~10.1.  It follows that, in any one of these gauges, the perturbation
theory will be renormalizable, in the sense that the divergences are removed
by a finite set of counterterms.  Furthermore, the analysis of Section~11.6
tells us that the only counterterms required are those that are symmetric
under the original global symmetry of the theory.  However, we should require
one further condition of our renormalization procedure: We should insist
that the counterterms preserve local gauge invariance, and, in particular,
preserve the property that $S$-matrix elements and the matrix elements of
gauge-invariant operators are independent of $\xi$.  This result was proved
to all orders in perturbation theory by 't Hooft and Veltman and by Lee and
Zinn-Justin.%
\footnote{G.~'t~Hooft and M.~J.~G.~Veltman, \emph{Nucl.\ Phys.\ B} \textbf{50},
318 (1972); B.~W.~Lee and J.~Zinn-Justin, \emph{Phys.\ Rev.\ D} \textbf{5},
3121, 3137, 3155 (1972), \textbf{D7}, 1049 (1973).}  Thus, in the gauge
defined by any finite value of $\xi$, we can, in principle, straightforwardly
compute a physical quantity to any order.  The gauges defined by the possible
values of $\xi$ are known as the \emph{renormalizability}, or $R_\xi$, gauges.

By taking the limit $\xi\to\infty$ of the $R_\xi$ gauges, we find a gauge
with very different simplifying features.  In this limit, the unphysical
degrees of freedom, which have masses proportional to $\sqrt{\xi}$,
disappear from the theory.  The gauge boson and Goldstone boson propagators
become:
\begin{equation}
  \frac{-i}{k^2 - m_A^2}\ \Bigl( g^{\mu\nu} - \frac{k^\mu k^\nu}{m_A^2}
    \Bigr),\qquad \frac{i}{k^2 - \xi m_A^2}\ \to 0 .
\end{equation}
The gauge boson propagator contains exactly the three spacelike polarization
states.  In this gauge, the only singularities of Feynman diagrams correspond
to the propagation of physical intermediate states.  Thus, the unitarity of
the $S$-matrix follows from the Cutkosky rules, as in the globally symmetric
theories considered in Section~7.3, without the need to worry about the
cancellation of unphysical states.%
\footnote{In the more sophisticated language of Section 16.4, the crucial
identity \eqref{eq:16.54}, which is required for the unitarity of the
$S$-matrix, is true manifestly.}  The $\xi\to\infty$ limit of the $R_\xi$
gauges thus gives the quantum-mechanical realization of the \emph{unitarity}
(or $U$) gauge, introduced in Eq.~\eqref{eq:20.12}.

It is not straightforward to prove renormalizability directly in the $U$
gauge.  In this gauge, the gauge boson propagator falls off more slowly than
$1/k^2$ at large $k$.  This signals trouble for the evaluation of loop
diagrams.  Typically, in fact, individual loop diagrams will diverge as
$\log\xi$ or worse as $\xi\to\infty$.  Still, the gauge invariance of the
$S$-matrix implies that these divergences must cancel in the sum of all
diagrams contributing to a given process, so that this sum has a smooth limit
as $\xi\to\infty$.  There is no difficulty of principle with the fact that we
use one gauge to prove the renormalizability of spontaneously broken gauge
theories and another gauge to prove their unitarity.  In fact, this method of
argument makes natural use of the underlying symmetries of the theory.

\subsection{Non-Abelian Analysis}

Now that we have thoroughly examined the $R_\xi$ gauges for an Abelian gauge
theory, we are ready to generalize to the non-Abelian case.  There is no
difficulty in being completely general, so let us consider a Yang-Mills gauge
theory with gauge group $G$, spontaneously broken by the vacuum expectation
value of a scalar field.

We will build on our classical analysis of this system following
Eq.~\eqref{eq:20.13}.  As in that analysis, it will be most convenient to
write the scalars as a multiplet $\phi_i$ of real-valued fields.  Then the
gauge transformation of the $\phi_i$ takes the form
\begin{equation}
  \delta\phi_i = -\alpha^a(x)\,T^a_{ij}\phi_j
  \label{eq:21.30}
\end{equation}
where the $T^a$ are real, antisymmetric representation matrices of $G$.
Similarly, the transformation of the gauge fields is
\begin{equation}
  \delta A^a_\mu = \partial_\mu\alpha^a - f^{abc}\alpha^b A^c_\mu
  = (D_\mu\alpha)^a .
  \label{eq:21.31}
\end{equation}
(If the gauge group is not simple, the coupling $g$ need not be the same for
every $a$.)  The Lagrangian invariant under these gauge transformations is
\begin{equation}
  \mathcal{L} = -\frac{1}{4}(F^a_{\mu\nu})^2 + \frac{1}{2}|D_\mu\phi|^2
    - V(\phi) ,
  \label{eq:21.32}
\end{equation}
with
\begin{equation}
  D_\mu\phi_i = \partial_\mu\phi_i + gA^{a\mu}\,T^a_{ij}\phi_j .
  \label{eq:21.33}
\end{equation}
Assume that the potential $V(\phi_i)$ is minimized at a point where some of
the components of $\phi$ acquire vacuum expectation values.  As in
\eqref{eq:20.16}, define
\begin{equation}
  \langle\phi_i\rangle = \phi_{0i} .
  \label{eq:21.34}
\end{equation}
We will expand $\phi_i$ about this value:
\begin{equation}
  \phi_i(x) = \phi_{0i} + \chi_i(x) .
  \label{eq:21.35}
\end{equation}
It will be convenient to divide the space of values $\chi_i$ into two
subspaces.  The vectors $T^a\phi_0$ correspond to symmetry transformations of
the vacuum expectation value of $\phi$.  The field fluctuations along these
directions are the Goldstone bosons.  Let $\{\eta_i^k\}$ be an orthonormal
basis for this subspace; then the unit vectors $\eta_i^k$ are in 1-to-1
correspondence with the Goldstone bosons.  The field fluctuations orthogonal
to all of the vectors $T^a\phi_0$ correspond to the (massive) physical scalar
fields of the spontaneously broken gauge theory.

In the discussion to follow, the vectors $T^a\phi_0$ will play an important
role.  We should then recall the notation for these vectors that we
introduced in Eq.~\eqref{eq:20.51}:
\begin{equation}
  F^{ai} = T^a_{ij}\phi_{0j} .
  \label{eq:21.36}
\end{equation}
The matrix $F^{ai}$ is not generally square; it has one row for each gauge
generator, and one column for each component of $\phi$.  However, many of
its elements are zero.  Its nonzero elements connect the spontaneously broken
gauge generators and the Goldstone bosons.  In Eq.~\eqref{eq:20.56}, we
showed that the gauge boson masses generated through the Higgs mechanism can
be written
\begin{equation}
  m^2_{ab} = g^2\,F^{aj}F^{bj} .
  \label{eq:21.37}
\end{equation}

To give a concrete example of a matrix $F^{aj}$, let us compute it in the
GWS electroweak theory.  Following the conventions introduced in
Eq.~\eqref{eq:20.14}, we should rewrite the Higgs field of the GWS model in
terms of four real scalar fields.  A convenient parametrization is
\begin{equation}
  \phi = \frac{1}{\sqrt{2}}\begin{pmatrix}
    \phi_1 + i\phi_2\\
    v + h + i\phi_3
  \end{pmatrix} .
  \label{eq:21.38}
\end{equation}
The fields $\phi^1, \phi^2, \phi^3$ are the Goldstone bosons, and $h$ is the
massive Higgs boson.  The vacuum state is simply
\begin{equation}
  \phi_0 = \frac{1}{\sqrt{2}}\begin{pmatrix}0\\ v\end{pmatrix} .
\end{equation}
The real representation matrices are
\begin{equation}
  T^a = -i\tau^a = -i\frac{\sigma^a}{2},\qquad
  T^Y = -iY = -i\frac{1}{2} .
\end{equation}
A simple computation then shows, for instance, that $T^1\phi_0$ equals
$v/2$ times a unit vector in the 1 direction.  Filling in the remaining
components of $F^{ai}$ with $a = 1, 2, 3, Y$ and $i = 1, 2, 3$, we find
\begin{equation}
  F^{ai} = \frac{v}{2}\ \begin{pmatrix}
    1 & 0 & 0\\
    0 & 1 & 0\\
    0 & 0 & 1\\
    0 & 0 & 0
  \end{pmatrix},
  \label{eq:21.39}
\end{equation}
We do not need to include the components of $F^{ai}$ along the direction of
the physical Higgs field $h$; the vectors $T^a\phi_0$ are all orthogonal to
this direction.

If we insert \eqref{eq:21.35} into \eqref{eq:21.32} as a change of
variables, we find, for the quadratic terms in the Lagrangian,
\begin{equation}
  \mathcal{L}_2 = -\frac{1}{4}A^{a\mu}\bigl( -g^{\mu\nu}\partial^2
    + \partial^\mu\partial^\nu \bigr)A^a_\nu + \frac{1}{2}(\partial_\mu\chi)^2
    + g\,\partial_\mu\chi_i\,A^{a\mu}F^{ai}
    + \frac{1}{2}(m^2_A)_{ab}A^a_\mu A^{b\mu} ,
  \label{eq:21.40}
\end{equation}
where $(m^2_A)_{ab}$ is the gauge boson mass matrix \eqref{eq:21.37} and
\begin{equation}
  M^2_{ij} = \frac{\partial^2V}{\partial\phi_i\partial\phi_j}
    \Big|_{\phi_0} .
  \label{eq:21.41}
\end{equation}
We proved in Eq.~\eqref{eq:11.13} that
\begin{equation}
  \eta_i^k M^2_{ij} = 0
  \label{eq:21.42}
\end{equation}
for all possible directions $\eta_i^k$ in the subspace spanned by the
$T^a\phi_0$, so the Goldstone bosons are massless.

To study the quantum theory of this system we start with the functional
integral
\begin{equation}
  Z = \int \mathcal{D}A\,\mathcal{D}\chi\ e^{i\int d^4x\ \mathcal{L}[A,\chi]} .
  \label{eq:21.43}
\end{equation}
Using the Faddeev-Popov gauge-fixing procedure, we define this integral,
analogously to \eqref{eq:21.10}, as
\begin{equation}
  Z = C'''\cdot \int \mathcal{D}A\,\mathcal{D}\chi\ \exp\Bigl[ i\int d^4x\
    \Bigl( \mathcal{L}[A,\chi] - \frac{1}{2\xi}(G)^2 \Bigr) \Bigr]\
    \det\Bigl( \frac{\delta G}{\delta\alpha} \Bigr) ,
  \label{eq:21.44}
\end{equation}
for an arbitrary gauge-fixing function $G(A,\chi)$.  The $R_\xi$ gauges are
defined by the choice
\begin{equation}
  G^a = \partial^\mu A^a_\mu - \xi g F^{ai}\chi_i .
  \label{eq:21.45}
\end{equation}
Note that $G$ involves only the components of $\chi$ that lie in the
subspace of the Goldstone bosons.

The gauge-fixing term adds to the Lagrangian the following set of quadratic
terms:
\begin{equation}
  \Bigl( -\frac{1}{2\xi}G^2 \Bigr) = \frac{1}{2\xi}A^{a\mu}\bigl(
    \partial^\mu\partial^\nu \bigr)A^a_\nu + g\,\partial^\mu A^{a\mu}F^{ai}\chi_i
    - \frac{\xi g^2}{2}\bigl[ F^{ai}\chi_i \bigr]^2 .
  \label{eq:21.46}
\end{equation}
The term that mixes $A^a_\mu$ and $\chi_i$ is arranged to cancel between
\eqref{eq:21.40} and \eqref{eq:21.46}.  The final quadratic Lagrangian for
the gauge and Goldstone boson fields is
\begin{equation}
  \mathcal{L}_2 = -\frac{1}{2}A^{a\mu}\Bigl[ \Bigl( -g_{\mu\nu}\partial^2
    + \Bigl( 1 - \frac{1}{\xi} \Bigr)\partial_\mu\partial_\nu
    - g^2F^{ai}F^{bi}g_{\mu\nu} \Bigr)A^{b\nu} \Bigr]
    - \frac{1}{2}\xi g^2(F^{ai}\chi_i)^2 .
  \label{eq:21.47}
\end{equation}
The mass matrices of gauge bosons and Goldstone bosons in this Lagrangian are
closely related to one another.  The gauge boson mass matrix is
\begin{equation}
  (m^2_A)_{ab} = g^2F^{ai}F^{bi} = g^2(FF^T)_{ab} .
  \label{eq:21.48}
\end{equation}
In an $R_\xi$ gauge, the timelike components of the gauge bosons acquire the
mass
\begin{equation}
  \xi m^2_A = \xi g^2(FF^T)_{ab} .
  \label{eq:21.49}
\end{equation}
At the same time, the Goldstone bosons acquire the mass matrix
\begin{equation}
  m^2_\varphi = \xi g^2(F^TF)_{ij} .
  \label{eq:21.50}
\end{equation}
The two matrices \eqref{eq:21.49} and \eqref{eq:21.50} have different numbers
of zero eigenvalues, but their nonzero eigenvalues are in 1-to-1
correspondence.  This is precisely the correspondence induced by the Higgs
mechanism between the massive gauge bosons and the Goldstone bosons that they
absorbed to gain mass.

Finally, we must construct the ghost Lagrangian.  This is found from the
gauge variation of the gauge-fixing term $G^a$.  Inserting \eqref{eq:21.30}
and \eqref{eq:21.31} into \eqref{eq:21.45}, we find
\begin{equation}
  \delta G^a = \Bigl( -\partial^2\delta^{ab} + \xi g^2\,(T^a\phi_0)\cdot
    T^b(\phi_0 + \chi) \Bigr)\alpha^b .
  \label{eq:21.51}
\end{equation}
Thus, the ghost Lagrangian is
\begin{equation}
  \mathcal{L}_{\mathrm{ghost}} = \bar{c}^a\Bigl[ -\partial^2\delta^{ab}
    - \xi g^2\,(T^a\phi_0)\cdot T^b(\phi_0 + \chi) \Bigr]c^b .
  \label{eq:21.52}
\end{equation}
Notice that the ghosts have exactly the same mass matrix \eqref{eq:21.49} as
the unphysical components of the gauge bosons.  This Lagrangian also contains
both the familiar coupling of the ghosts to the gauge fields and the coupling
to the physical Higgs fields that we found in the Abelian case \eqref{eq:21.15}.

We have now computed the kinetic energy terms for gauge fields, scalar
fields, and ghosts in an $R_\xi$ gauge.  It is straightforward to convert
these results to the calculation of propagators for these fields; the
computations are exactly the same as in the Abelian case.  We find for the
three propagators the standard forms, with the gauge and Goldstone boson
masses given by the matrices \eqref{eq:21.48}--\eqref{eq:21.50}.  This
completes the quantization of spontaneously broken gauge theories in the
$R_\xi$ gauges.

\section{The Goldstone Boson Equivalence Theorem}

The quantization of spontaneously broken gauge theories in the $R_\xi$ gauges
leads to a remarkable and important result, the \emph{Goldstone boson
equivalence theorem}: at high energies, the amplitudes for emission or
absorption of longitudinally polarized gauge bosons become equal to the
amplitudes for the corresponding Goldstone bosons.  This theorem has
important applications to the physics of the weak interactions, since the
longitudinal components of the $W$ and $Z$ bosons dominate their production
and decay at high energies.

To prepare for a discussion of longitudinal vector bosons, we need some
simple kinematics.  A vector boson at rest has momentum
$k^\mu = (m, 0, 0, 0)$ and a polarization vector that is a linear combination
of the three orthogonal unit vectors
\begin{equation}
  (0, 1, 0, 0),\qquad (0, 0, 1, 0),\qquad (0, 0, 0, 1) .
  \label{eq:21.57}
\end{equation}
If we boost this particle along the 3 axis, its momentum boosts to
\begin{equation}
  k^\mu = (E_k, 0, 0, k) .
\end{equation}
The three possible polarization vectors are now the three unit vectors
satisfying
\begin{equation}
  \epsilon^2 = 0,\qquad \epsilon^{*}\cdot\epsilon = -1 .
  \label{eq:21.58}
\end{equation}
Two of these are the first two vectors in \eqref{eq:21.57}; these give the
transverse polarizations.  The third vector satisfying \eqref{eq:21.58} is
the longitudinal polarization vector
\begin{equation}
  \epsilon^\mu_L(k) = \Bigl( \frac{k}{m}, 0, 0, \frac{E_k}{m} \Bigr) ,
  \label{eq:21.59}
\end{equation}
which is the boost of the third vector in \eqref{eq:21.57}.  An important
and somewhat counterintuitive feature of \eqref{eq:21.59} is that it becomes
increasingly parallel to $k^\mu$ as $k$ becomes large.  In fact, component by
component,
\begin{equation}
  \epsilon^\mu_L(k) = \frac{k^\mu}{m} + \mathcal{O}(m/E_k)
  \label{eq:21.60}
\end{equation}
as $k\to\infty$.  Since the components of $k^\mu$ are growing as $k$, this
statement is consistent with the requirement that $\epsilon^{*}\cdot k = 0$
while $k\cdot k = m^2$.

With this kinematics situation in mind, let us analyze the Ward identity
satisfied by a gauge current matrix element between on-shell states.  It is
simplest to work in Lorentz gauge ($\xi = 0$), where the gauge-fixing term
\eqref{eq:21.45} does not involve the Goldstone boson fields.  The Ward
identity can then be written as follows:
\begin{equation}
  k_\mu\langle \beta|\ J^{\mu a}(k)\ |\alpha\rangle = 0 ,
  \label{eq:21.64}
\end{equation}
between on-shell states.  Writing the matrix element as the sum of two
pieces---first, the current can couple directly into a one-particle-irreducible
vertex function $\Gamma^\mu(k)$; second, for a spontaneously broken gauge
theory, there is an additional term, not one-particle-irreducible, in which
the current creates a Goldstone boson and it is this particle that couples to
the external states through a 1PI vertex $\Gamma(k)$---we find
\begin{equation}
  k_\mu\Gamma^\mu(k) + k_\mu\,(igF)\ \frac{1}{k^2}\ \Gamma(k) = 0 .
  \label{eq:21.65}
\end{equation}
Let us write the relation linking the gauge current and the Goldstone boson
state as
\begin{equation}
  \langle 0|J^\mu|\pi(k)\rangle = -iF k^\mu ,
  \label{eq:21.62}
\end{equation}
as in Eq.~\eqref{eq:20.46}.  Then the argument leading to Eq.~\eqref{eq:20.56}
tells us that the gauge boson mass is given by
\begin{equation}
  m = gF ,
  \label{eq:21.63}
\end{equation}
where $g$ is the gauge boson coupling constant.  With these identifications,
we can write the Ward identity \eqref{eq:21.64} as
\begin{equation}
  k_\mu\Gamma^\mu(k) = m\,\Gamma(k) .
  \label{eq:21.66}
\end{equation}
Now use this equation in the limit of large gauge boson momentum.  Since the
gauge boson vertex is one-particle-irreducible, the momenta of propagators
inside the vertex are not, in general, collinear with $k^\mu$.  Then,
according to \eqref{eq:21.60}, we may replace $k^\mu/m$ by the longitudinal
polarization vector.  Notice that this would not be permissible (but, also,
is not necessary) in the second term of \eqref{eq:21.65}.  Our final result
is
\begin{equation}
  \epsilon^\mu_L(k)\,\Gamma_\mu(k) = \Gamma(k) ,
  \label{eq:21.67}
\end{equation}
as $k\to\infty$, with an error of order $m^2/k^2$.  That is, in the
high-energy limit, the couplings of longitudinal gauge bosons become
precisely those of their associated Goldstone bosons.

The equivalence theorem can be derived in another way, using the counting of
physical states in spontaneously broken gauge theories, which we discussed
below Eq.~(21.26).  In the previous section, we saw that, at least
at the tree level, unitarity is maintained in spontaneously broken gauge
theories by the cancellation of diagrams that produce timelike-polarized
gauge bosons against diagrams that produce Goldstone bosons.

The situation is most clear in Feynman-'t Hooft gauge.  There, the numerator
of the gauge boson propagator is $-g^{\mu\nu}$.  We can write this in terms
of polarization vectors as
\begin{equation}
  -g^{\mu\nu} = \sum_{\epsilon} \epsilon^{\mu*}(k)\epsilon^\nu(k)
    - \frac{k^\mu k^\nu}{m^2} .
  \label{eq:21.68}
\end{equation}
The last term is the contribution from unphysical timelike polarization
states.  The unitarity of the $S$-matrix requires that, when a Cutkosky cut
through a diagram puts a gauge boson propagator on-shell, the contribution of
this piece must be canceled by a Cutkosky cut that runs through a Goldstone
boson line.  The required cancellation is
\begin{equation}
  \Bigl| \frac{k^\mu}{m}\Gamma_\mu(k) \Bigr|^2 + |\Gamma(k)|^2 = 0 ,
  \label{eq:21.69}
\end{equation}
or, diagrammatically,
\begin{equation}
  \text{(timelike gauge boson cut)} + \text{(Goldstone boson cut)} = 0 .
\end{equation}
Once again, since $\Gamma^\mu(k)$ is a one-particle-irreducible vertex, we
can use \eqref{eq:21.60} to replace $(k^\mu/m)$ by the longitudinal
polarization vector $\epsilon^\mu_L(k)$ for a high energy gauge boson.  Then
\eqref{eq:21.69} becomes just the square of \eqref{eq:21.67}.

Through these formal arguments, we can see, at least to the tree level in
processes with single gauge boson emission, that the equivalence theorem must
be valid.  However, it is much more illuminating to see the equivalence
theorem at work in explicit calculations for interesting physical processes.
We will now illustrate its influence in two examples.%
\footnote{For a general proof of the equivalence theorem, including the case
of multiple absorptions and emissions of massive vector bosons, see
M.~S.~Chanowitz and M.~K.~Gaillard, \emph{Nucl.\ Phys.\ B} \textbf{261}, 379
(1985).}

\subsection{Top Quark Decay}

The first example is the weak decay of the top quark.  This charge $+2/3$
quark is sufficiently heavy that it can decay to a real $W^+$ through
$t\to W^+ + b$.  The diagram for this decay is given by the simple gauge
vertex shown in Fig.~\ref{fig:21.4}.

\begin{figure}[t]
  \centering
  \includegraphics[width=0.45\textwidth]{fig21.4.png}
  \caption{Decay of a $t$ quark into $W^+ + b$.}
  \label{fig:21.4}
\end{figure}

Let us first try to guess the magnitude of the top quark width.  The squared
matrix element will contain a factor of $g^2$, times some expression with
dimensions of mass.  Since the width should be large if the top quark mass is
heavy, a first guess might be
\begin{equation}
  \Gamma \sim \frac{g^2}{4\pi}\,m_t .
  \label{eq:21.70}
\end{equation}
The correct expression, however, turns out to be enhanced by a factor of
$(m_t/m_W)^2$.

The amplitude for this decay can be read from Eq.~\eqref{eq:20.80}:
\begin{equation}
  i\mathcal{M} = \frac{ig}{\sqrt{2}}\ \bar{u}(q)\gamma^\mu\Bigl(
    \frac{1-\gamma^5}{2} \Bigr)u(p)\,\epsilon^{*}_\mu(k) .
  \label{eq:21.71}
\end{equation}
(We set the relevant CKM factor equal to 1.)  We will now turn this
amplitude into an expression for the decay rate of the top quark.  For
simplicity, we will ignore the mass of the $b$ quark in this computation.

Squaring the amplitude in \eqref{eq:21.71} according to our standard
methods, and then averaging over initial and summing over final spins, we
find
\begin{equation}
  \sum_{\text{spins}}|\mathcal{M}|^2 = \frac{g^2}{2}\ \mathrm{tr}\Bigl[
    \slashed{q}\gamma^\mu\slashed{p}\gamma^\nu\ \frac{1-\gamma^5}{2}
    \Bigr]\ \sum_{\text{polarizations}}\ \epsilon^{*}_\mu(k)\epsilon_\nu(k)
  .
  \label{eq:21.72}
\end{equation}
We can sum explicitly over physical gauge boson polarizations by inserting
the expression for the polarization sum.  This gives
\begin{equation}
  \sum_{\text{spins}}|\mathcal{M}|^2 = \frac{g^2}{2}\ \Bigl[
    4\Bigl( -g^{\mu\nu} + \frac{k^\mu k^\nu}{m_W^2} \Bigr)\ \cdots
    \Bigr] ,
  \label{eq:21.73}
\end{equation}
where the contribution of the longitudinal polarization, proportional to
$k^\mu k^\nu/m_W^2$, gives the dominant term.  Evaluating the traces, one
finds that the longitudinal contribution is enhanced by a factor of
$(m_t/m_W)^2$, precisely as required by the equivalence theorem.  The decay
width of the top quark is
\begin{equation}
  \Gamma(t\to W^+b) = \frac{g^2}{64\pi}\ \frac{m_t^3}{m_W^2}\ \cdots ,
  \label{eq:21.74}
\end{equation}
which is much larger than the naive estimate \eqref{eq:21.70}.  This is a
direct illustration of the Goldstone boson equivalence theorem: the
longitudinal $W$ boson in the final state behaves like the corresponding
Goldstone boson, and the amplitude is enhanced by the large $t$-quark mass.

\subsection{$e^{+}e^{-}\to W^{+}W^{-}$}

The second example is the process $e^{+}e^{-}\to W^{+}W^{-}$ in the
weak-interaction gauge theory.  The relevant diagrams are shown in
Fig.~\ref{fig:21.7}.  This process is interesting because the cross section
for production of two longitudinal $W$ bosons must, by the equivalence
theorem, equal the cross section for production of two Goldstone bosons at
high energy.

\begin{figure}[t]
  \centering
  \includegraphics[width=0.6\textwidth]{fig21.7.png}
  \caption{Diagrams contributing to $e^{+}e^{-}\to W^{+}W^{-}$ in the weak
    interaction gauge theory.}
  \label{fig:21.7}
\end{figure}

Let us first compute the amplitude for production of a pair of charged
scalars (the Goldstone bosons).  Assembling vertices from
Fig.~\ref{fig:21.8},

\begin{figure}[t]
  \centering
  \includegraphics[width=0.6\textwidth]{fig21.8.png}
  \caption{Feynman rules of the weak-interaction gauge theory for electrons
    and scalars coupling to photons and $Z$ bosons.}
  \label{fig:21.8}
\end{figure}

we find that, for an electron of either helicity, the
amplitude to annihilate to scalars through a virtual photon is
\begin{equation}
  i\mathcal{M}(e\bar{e}\to\gamma^{*}\to\phi^+\phi^-)
  = ie^2\ \bar{v}(p')\gamma^\mu u(p)\ \frac{1}{q^2}\ (k_+ - k_-)_\mu ,
  \label{eq:21.92}
\end{equation}
where $k_+$, $k_-$ are the momenta of the scalars and $q = k_+ + k_-$.  The
corresponding amplitude for annihilation through a virtual $Z^0$ depends on
the $e^{+}e^{-}$ helicities.  Adding these contributions to the preceding
expression, we find
\begin{equation}
  \begin{aligned}
    i\mathcal{M}(e^-_Le^+_R \to \phi^+\phi^-) &= ie^2\ \bar{v}_R\gamma^\mu
      u_L\ \Bigl[ \frac{1}{q^2} + \frac{\bigl( \tfrac{1}{2} - \sin^2\theta_W
        \bigr)^2}{\sin^2\theta_W\cos^2\theta_W}\ \frac{1}{q^2 - m_Z^2}
      \Bigr](k_+ - k_-)_\mu ,\\
    i\mathcal{M}(e^-_Re^+_L \to \phi^+\phi^-) &= ie^2\ \bar{v}_L\gamma^\mu
      u_R\ \Bigl[ \frac{1}{q^2} + \frac{(\sin^2\theta_W)^2}{\sin^2\theta_W
        \cos^2\theta_W}\ \frac{1}{q^2 - m_Z^2} \Bigr](k_+ - k_-)_\mu .
  \end{aligned}
  \label{eq:21.93}
\end{equation}
Notice that, in the high-energy limit, the amplitude for the annihilation of
right-handed electrons cancels down to
\begin{equation}
  i\mathcal{M}(e^-_Re^+_L \to \phi^+\phi^-)
  = ie^2\ \bar{v}_L\gamma^\mu u_R\ \frac{1}{q^2}\,(k_+ - k_-)_\mu ,
  \label{eq:21.94}
\end{equation}
which is just the amplitude for an $e_R$, with $Y = -1$, to couple to a
$\phi^+$, with $Y = 1/2$, through the $U(1)$ gauge boson $B_\mu$ with
coupling constant $g' = e/\cos\theta_W$.  This expression reflects the fact
that the $\phi^+$ has no direct coupling to the $SU(2)$ gauge bosons.
Similarly, the amplitude for left-handed electrons tends to
\begin{equation}
  i\mathcal{M}(e^-_Le^+_R \to \phi^+\phi^-)
  = ie^2\ \bar{v}_R\gamma^\mu u_L\ \frac{1}{4\cos^2\theta_W\sin^2\theta_W}\
    \frac{1}{q^2}\,(k_+ - k_-)_\mu ,
  \label{eq:21.95}
\end{equation}
in the high-energy limit.  This has the structure of a coherent sum of
amplitudes with $B^\mu$ and $A^\mu$ exchange.  In just the way that we saw in
Chapter~11, the symmetry structure of a gauge theory with spontaneously
broken symmetry is recovered in the high-energy limit.

Now let us compare these results to a direct calculation of the $W^{+}W^{-}$
production amplitude in the weak interaction gauge theory.  Begin with the
case of an initial $e^-_R$.  Since the coupling of the electron to the $W^-$
is purely left-handed, the third diagram of Fig.~\ref{fig:21.7} vanishes in
this case, so the computation is a bit easier.  The first two diagrams of
Fig.~\ref{fig:21.7} have exactly the same structure and sum to
\begin{equation}
  i\mathcal{M}(e^-_Re^+_L \to W^+W^-) = \bar{v}_R\gamma^\mu u_R\ \frac{
    ie^2}{\sin^2\theta_W}\ \Bigl[ \cdots \Bigr]\
    \epsilon^{*}_\mu(k_+)\epsilon^{*}_\nu(k_-) ,
  \label{eq:21.96}
\end{equation}
where the expression in brackets contains the photon and $Z^0$ propagators
and the $WW\gamma$ and $WWZ$ vertices.  This equation is valid in any of
the $R_\xi$ gauges, since, if we ignore the electron mass,
\begin{equation}
  q^\mu\ \bar{v}_R\gamma_\mu u_R = 0 .
  \label{eq:21.97}
\end{equation}
The second line of Eq.~\eqref{eq:21.96} contains the enhancement for
longitudinal $W$ bosons mentioned above.  If we approximate the longitudinal
polarization vectors by \eqref{eq:21.60} and drop terms that do not grow as
$s\to\infty$, this line becomes proportional to
\begin{equation}
  \frac{1}{m_W^2}\ (k_+ - k_-)^\lambda\ \cdots ,
  \label{eq:21.98}
\end{equation}
which is precisely the structure of the Goldstone boson amplitude
\eqref{eq:21.93}.  Carrying out the algebra, one confirms that the amplitude
for producing two longitudinal $W$ bosons equals the amplitude for producing
two Goldstone bosons, in accord with the equivalence theorem.

The Goldstone boson equivalence theorem has important physical consequences.
At energies much larger than $m_W$, the production of longitudinal $W$ and
$Z$ bosons is enhanced relative to the naive expectation, and the cross
sections are governed by the Goldstone boson interactions.  This means, for
example, that the cross section for $e^{+}e^{-}\to W^{+}W^{-}$ grows with
energy in a way that is determined by the structure of the Higgs sector.  At
very high energies, the theory would violate unitarity unless the $W$ and $Z$
boson interactions satisfy the relations imposed by the equivalence theorem.
This requirement was one of the original motivations for the Higgs
mechanism.

\section{One-Loop Corrections to the Weak-interaction Gauge Theory}

In Section~11.6, we presented a general argument that, in a theory with
spontaneously broken symmetry, relations between parameters and observables
that hold at the classical level can be modified by loop corrections.  These
corrections will be finite, and will in fact be definite predictions from the
quantum structure of the field theory.  Though we discussed this idea
originally in theories with spontaneously broken global symmetry, it applies
equally well to theories with spontaneously broken gauge symmetry.  In this
section, we will apply this idea to derive finite one-loop corrections to
relations in the weak-interaction gauge theory.

It is easy to find zeroth-order natural relations in the electroweak theory.
The leading-order predictions given in Section~20.2 involve a relatively
small number of free parameters.  Many of these predictions are made for
energies at which the quark and lepton masses can be ignored; then they
depend only on the coupling constants $g$ and $g'$ and the vacuum
expectation value $v$, which sets the scale of spontaneous symmetry breaking.
The remaining ingredients of the weak-interaction theory are given in terms
of these parameters; for example,
\begin{equation}
  m_W = g\frac{v}{2},\qquad
  m_Z = \sqrt{g^2 + g'^2}\,\frac{v}{2},\qquad
  e = \frac{gg'}{\sqrt{g^2+g'^2}},\qquad
  G_F = \frac{1}{\sqrt{2}}\ \frac{g^2}{m_W^2} .
  \label{eq:21.109}
\end{equation}
Even in this set of quantities, we have four relations that depend on three
underlying parameters, so there is one relation of observable quantities
that is independent of the parameters of the Lagrangian.

Since many of the predictions of the weak interaction gauge theory are
determined by the parameter $\sin^2\theta_W$, it is useful to define
$\sin^2\theta_W$ in terms of observables and then use this definition as a
basis for constructing natural relations.  In our discussion of the precision
tests of electroweak theory in Section~20.2, we used the definition
\begin{equation}
  \sin^2\theta^s_W = 1 - \frac{m_W^2}{m_Z^2} .
  \label{eq:21.110}
\end{equation}
as a standard for comparison of different experiments.  But since the three
most accurately known weak-interaction observables are $\alpha$, $G_F$, and
$m_Z$, it is useful to construct another physical definition of
$\sin^2\theta_W$ based on these three quantities.  Define $\theta_0$ such
that
\begin{equation}
  \sin^2\theta_0\cos^2\theta_0 = \frac{\pi\alpha^{*}}{\sqrt{2}G_Fm_Z^2} ,
  \label{eq:21.111}
\end{equation}
where $\alpha^{*}$ is the running coupling constant of QED evaluated at the
scale $Q^2 = m_Z^2$.  The renormalization group insists that it is the value
of the electric charge at the weak-interaction scale that enters precision
electroweak predictions, and this observation is confirmed by summing
radiative correction diagrams involving light quarks and leptons.  The
current best values of the quantities in Eq.~\eqref{eq:21.111} give
\begin{equation}
  s^2_0 = \sin^2\theta_0 = 0.2307 \pm 0.0005 .
  \label{eq:21.112}
\end{equation}
Thus, this quantity provides a very accurate standard of reference.

Once Eq.~\eqref{eq:21.111} is taken to define a reference value of
$\sin^2\theta_W$, the equations of Section~20.2 that connect
$\sin^2\theta_W$ to other observables become zeroth-order natural relations.
For example, the tree-level equations
\begin{equation}
  \frac{m_W^2}{m_Z^2} = \cos^2\theta_W,\qquad
  \mathcal{A}_{LR} = \frac{\bigl( \tfrac{1}{2} - \sin^2\theta_W \bigr)^2
    - \bigl( \sin^2\theta_W \bigr)^2}
    {\bigl( \tfrac{1}{2} - \sin^2\theta_W \bigr)^2
    + \bigl( \sin^2\theta_W \bigr)^2} ,
  \label{eq:21.113}
\end{equation}
are natural relations linking four observables of the weak interactions.
The corrections to these relations will be well-defined predictions of the
theory.

In principle, we could now compute all of the one-loop diagrams that correct
the parameters $m_Z$, $m_W$, $G_F$, $\alpha$, and $\mathcal{A}_{LR}$.
However, this is a very complicated exercise, requiring an extensive
technical apparatus.%
\footnote{A detailed theoretical discussion of one-loop corrections to the
electroweak theory can be found in W.~Hollik, \emph{Fortschr.\ d.\ Physik}
\textbf{38}, 165 (1990).}  In this section we will focus on radiative
corrections from one simple source that can be considered independently.
Aside from the question of anomalies, the electroweak theory does not
restrict the number of quark or lepton generations.  Thus, it is sensible,
and gauge invariant, to compute the one-loop corrections due to one quark or
lepton doublet.  For definiteness, we consider the effects of the $(t, b)$
quark doublet.

By focusing on the radiative corrections due to heavy quarks, we
dramatically simplify the computational task before us.  The various
observables of the weak-interaction gauge theory are extracted from the
measurement of scattering amplitudes with light fermions, leptons or quarks,
in the initial and final states.  For example, $G_F$ is measured from the
strength of a low-energy weak-interaction process, usually chosen to be the
rate of muon decay: $\mu^{-}\to e^{-}\bar{\nu}_e\nu_\mu$.  For any such
process, there are one-loop corrections of many kinds, as shown in
Fig.~21.11.  In addition to corrections to the vector boson
propagator, there are vertex corrections, box diagrams, and diagrams with
real photon emission.  In general, the contributions of the various classes
of diagrams are not gauge invariant; rather, gauge invariance results from
cancellations between the classes of diagrams in Fig.~21.11(b),
(c), and (d).  However, since heavy quarks do not couple directly to the
light leptons, the $(t, b)$ doublet contributes only the single diagram shown
in Fig.~21.11(f), which must be gauge invariant by itself.  This
same conclusion applies to the $(t, b)$ correction to other leptonic weak
interaction processes.  If we ignore the CKM angles that mix the $t$ and $b$
with other species, the conclusion extends also to weak-interaction
processes involving light quarks.

A similar situation occurs with other species of particles, such as those of
the Higgs sector.  The coupling of Higgs sector particles to a light quark or
lepton is proportional to the fermion's mass, which we can often ignore.
Thus the most important contributions from Higgs-sector particles are
propagator corrections.  The case in which the spontaneous symmetry breaking
is produced by a single scalar field $\phi$ is particularly straightforward
to analyze; this is done in Problem~21.4.  Loop corrections from particles
that do not couple directly to the external fermions are often termed
\emph{oblique}, since they enter the low-energy weak interactions only
indirectly.

\subsection{Influence of Heavy Quark Corrections}

Our task, then, is to compute the corrections to relations \eqref{eq:21.113}
due to the $(t, b)$ doublet.  These two relations depend on five observable
quantities---$m_Z$, $m_W$, $\mathcal{A}_{LR}$, $\alpha$, and $G_F$---with
the last two parameters entering through $\theta_W$ and Eq.~\eqref{eq:21.111}.
We will express these five quantities as functions of the bare parameters
$g$, $g'$, and $v$, with corrections proportional to combinations of $t$ and
$b$ vacuum polarization diagrams.  The zeroth-order terms will naturally
cancel out when we compute the corrections to the relations \eqref{eq:21.113}.

The loop amplitudes that we require are the one-loop vacuum polarization
diagrams of the $t$ and $b$ quarks.  To
deal with these contributions most straightforwardly, we introduce a uniform
notation for vacuum polarization amplitudes.  Denote the vacuum
polarization amplitude involving the gauge bosons $I$ and $J$ as
\begin{equation}
  i\Pi^{\mu\nu}_{IJ}(q) ,
  \label{eq:21.114}
\end{equation}
where $I$ and $J$ may be $\gamma$, $W$, or $Z$.  When the gauge bosons are
massive, the vacuum polarization amplitudes need not be transverse by
themselves, so $\Pi^{\mu\nu}_{IJ}(q)$ need not vanish at $q^2 = 0$.  Thus, we
will change our notation from the case of QED and write the decomposition of
$\Pi^{\mu\nu}_{IJ}(q)$ into tensor structures as
\begin{equation}
  \Pi^{\mu\nu}_{IJ}(q) = \Pi_{IJ}(q^2)\,g^{\mu\nu} - \Delta_{IJ}(q^2)\,q^\mu q^\nu
  .
  \label{eq:21.115}
\end{equation}
In all of the examples to follow, the factors $q^\nu$ will dot into currents
of light leptons, to give zero as in Eq.~\eqref{eq:21.97}.  Thus the form
factor $\Delta_{IJ}(q^2)$ will drop out of our calculations.  Our previous
result that $\Pi^{\mu\nu}(q)$ vanishes in QED at $q^2 = 0$ appears in this
formalism as the set of constraints
\begin{equation}
  \Pi_{\gamma\gamma}(0) = \Pi_{\gamma Z}(0) = 0 .
  \label{eq:21.116}
\end{equation}
For the other amplitudes, our sign conventions are chosen so that a positive
value of $\Pi_{IJ}(m^2)$ gives a positive mass shift to the gauge boson.  Let
us also define
\begin{equation}
  \Pi'_{IJ} = \frac{d\Pi_{IJ}}{dq^2}\Big|_{q^2=0},
  \label{eq:21.117}
\end{equation}
this is the quantity we called $\Pi(0)$ in Eq.~\eqref{eq:7.73}.

Now we use this notation to write the loop corrections to each of the five
observables.  The first two diagrams in Fig.~21.12 are simply mass
corrections, and so, straightforwardly,
\begin{equation}
  m_Z^2 = (g^2 + g'^2)\frac{v^2}{4} + \Pi_{ZZ}(m_Z^2),\qquad
  m_W^2 = g^2\frac{v^2}{4} + \Pi_{WW}(m_W^2) .
  \label{eq:21.118}
\end{equation}
Note that both vacuum polarization amplitudes are evaluated at the poles in
the respective propagators.  To evaluate the shift of $\alpha$ by one-loop
corrections, we consider the effect of Fig.~21.12(c) on the
low-energy Coulomb potential.  The values of the leading-order propagator and
the one-loop correction combine to give the factors
\begin{equation}
  \frac{e^2}{q^2}\ \Bigl( 1 + \Pi'_{\gamma\gamma}(q^2)\ \cdots \Bigr) ,
  \label{eq:21.119}
\end{equation}
where, in this equation, $e^2$ is given in terms of bare variables as in
\eqref{eq:21.109}.  Thus, the observed value of $\alpha$, in the limit
$q^2\to 0$, is modified according to
\begin{equation}
  \alpha = \alpha_0\ \Bigl( 1 - \Pi'_{\gamma\gamma}(0)\ \cdots \Bigr) .
  \label{eq:21.120}
\end{equation}
Similarly, the corrections to $G_F$ and $\mathcal{A}_{LR}$ can be expressed
in terms of the vacuum polarization amplitudes of the $(t,b)$ doublet.

The contributions of the $(t, b)$ doublet to the vacuum polarization
amplitudes can be computed explicitly.  The diagrams are the standard
one-loop vacuum polarization diagrams with $t$ and $b$ quarks in the loop.
The result is that the corrections to the natural relations depend on the
difference of the squared masses of the $t$ and $b$ quarks,
\begin{equation}
  \Delta\rho \propto \frac{m_t^2 - m_b^2}{m_W^2} .
  \label{eq:21.121}
\end{equation}
This is the famous result that the $t$ quark mass dependence of the
electroweak radiative corrections is quadratic, not logarithmic.  The
quantity $\Delta\rho$ modifies the relation between $m_W$ and $m_Z$, so that
\begin{equation}
  \frac{m_W^2}{m_Z^2} = \cos^2\theta_W\ \Bigl( 1 - \frac{\Delta\rho}
    {\cos^2\theta_W} \Bigr) .
  \label{eq:21.122}
\end{equation}
The dependence of the weak-interaction observables on $m_t^2$ is a distinctive
prediction of the theory, and it was used to estimate the top quark mass
before its discovery.

\subsection{The Effect of $m_t$}

The one-loop corrections to the weak-interaction observables are dominated by
the quadratic dependence on the top quark mass.  This dependence arises from
the fact that the $(t, b)$ doublet is split, with $m_t \gg m_b$.  The
isospin-breaking effect of the large $t$ mass modifies the relative strength
of the neutral and charged current interactions.

The result of the full calculation is that the natural relation
\eqref{eq:21.113} is corrected by a well-defined amount, which depends on
$m_t$ and on the properties of the Higgs sector.  The comparison of the
predicted corrections with the experimental values of the weak-interaction
observables provides a stringent test of the theory, and the agreement
between theory and experiment is one of the great successes of the standard
model.  In particular, the observed value of the $\rho$ parameter is
consistent with a top quark mass in the range that has been measured
directly, and with a light Higgs boson.  This agreement gives confidence in
the validity of the standard model at the quantum level.

The radiative corrections to the weak-interaction gauge theory also provide
a powerful way to constrain the properties of new particles beyond the
standard model.  Any new particle that couples to the $W$ and $Z$ bosons
will contribute to the vacuum polarization amplitudes, and thus to the
corrections to the natural relations.  The precise agreement between theory
and experiment for the measured observables places strong constraints on the
masses and couplings of such new particles.  This method has been used, for
example, to constrain the mass of the Higgs boson and the possible existence
of additional $Z'$ bosons.

\problems

\textbf{21.1} \emph{The Goldstone boson equivalence theorem in the Abelian
model.}  In this problem, verify the Goldstone boson equivalence theorem in
the Abelian spontaneously broken gauge theory of Section~21.1.  Compute the
amplitude for fermion-antifermion annihilation to produce a longitudinal
gauge boson, and show that it equals the amplitude for producing the
Goldstone boson at high energy.

\textbf{21.2} \emph{$\xi$ independence of $S$-matrix elements.}  In the
Abelian model of Section~21.1, verify explicitly that the tree-level
fermion-fermion scattering amplitude is independent of the gauge parameter
$\xi$, by summing the diagrams of Fig.~\ref{fig:21.2}.

\textbf{21.3} \emph{The $W$ boson propagator in the $R_\xi$ gauges.}  Verify
the form of the gauge boson propagator \eqref{eq:21.17} by inverting the
matrix \eqref{eq:21.16}.  Show that the unphysical poles in the propagator
cancel in $S$-matrix elements.

\textbf{21.4} \emph{One-loop corrections from the Higgs sector.}  In the
Abelian spontaneously broken gauge theory, compute the one-loop corrections
to the gauge boson mass and to the natural relation between the gauge boson
mass and the coupling constants, coming from the Higgs boson loop.  Show
that these corrections are finite and well-defined.
